\section{Applications: transfers and recorded obstructions}
\label{sec:apps}

The fourth stage of the plan applies admitted Round32 and Round33 equations
to related problems and records, for each equation it applies, either a
labelled transfer to a named model with its own exact checks or the
obstruction that prevents the transfer (Section~\ref{sec:plan}). The BD1 and
BD2 contracts were drafted during the BC production, since no BD item
depends on a BC outcome, and were frozen together on 25 September 2026 at
05:09:10 UTC, after the skeptic's pre-freeze review of the drafts (34 blocking
and 21 non-blocking edits applied). The review found the drafted BD2 width
target of $1/10^{10}$ non-discriminating, and the advisor tightened it to
$1/10^{16}$ at the cutoff $D=8$ before the freeze. BD1 applies the Round32
centre-symmetry equations (the AW1 link-flip lemma, parity theorem and
first-order Wilson coefficient) to other gauge groups. BD2 applies admitted
SU(2) equations to new observables, to a finite graph with independent
couplings and to a new dimension. Both gates were recorded at 07:52 and 07:53
UTC on 25 September; with them all eight Round33 investigations are gated.
Table~\ref{tab:apps-summary} at the end of the section summarizes what
transferred and what did not.

\subsection{BD1: centre symmetry across gauge groups}
\label{sec:bd1}

\textbf{Contribution \projecttag{HNM-C-BD1}.} The BD1 contract, recorded
under the title ``Applications I: centre-symmetry transfer of the parity
theorem, the link-flip lemma and the first-order Wilson mean across gauge
groups, with the SU(3) and SO(3) obstructions and the SU(2) area-parity
corollary'', asks which of three AW1 statements survive when SU(2) is
replaced by SU(3), SU(4), SU(5), U(1), Z2 or SO(3), and asks for an exact
counterexample wherever one does not
(\repo{research/round33/advisor/selection-bd1.md}{selection note}). It is
investigation~7 of~8 and paired. The forward producer computes by
characters: Peter--Weyl, the Pieri rule, Clebsch--Gordan series and charge
counting. The reverse producer, under premise isolation, computes by Weyl
integration: constant terms on the maximal torus. Both received
\texttt{AGENTS.md}, the contract and 21 shared premises. The target is exact
agreement of the two routes on every moment and coefficient, and nonzero
SU(3) and SO(3) obstruction coefficients; the contract notes that the
obstruction cells are discriminating, since a zero value would falsify
them. It lists 22 controls. Its gate is recorded under the title ``Hruday
centre-symmetry transfer across gauge groups with recorded obstructions''.

\subsubsection{Reviewed scope}
\label{sec:bd1-scope}
The \repo{research/round33/advisor/bd1-gate.json}{BD1 gate} records the
verdict \recid{accepted_within_scope}, with sub-label
\recid{transfer_to_named_model} and secondary sub-label
\recid{obstruction_recorded}. Its gate fields set
\recid{flip_transfer_claimed}, \recid{parity_transfer_claimed},
\recid{obstructions_recorded}, \recid{area_parity_limit_claimed} and
\recid{model_is_finite_graph} to true, and \recid{transfers_to_aq},
\recid{uniqueness_of_ground_state_claimed}, \recid{rate_in_a_claimed},
\recid{continuum_claim}, \recid{weak_coupling_claim} and
\recid{scientific_priority_verified} to false. The flip scope names SU(2),
SU(4), U(1) and Z2 on their one-plaquette and box models, with no AM2, AV1 or
AQ statement for any group other than SU(2). The accepted statement, quoted
verbatim below in the gate's ASCII notation, is the scope of every claim in
this subsection.
\begin{gatequote}
Under the frozen BD1 convention (per-link electric term 8 C\_2 with
C\_F=(N\textasciicircum{}2-1)/(2N) for SU(N), C\_2(n)=n\textasciicircum{}2 for
U(1), C\_2(l)=l(l+1) for SO(3) and C\_2=1 on the odd Z2 state; face energy 32
C\_F; magnetic term -(tau/3) W with W = Re chi\_fund/dim fund; creation sign
psi=e\textasciicircum{}\{-C\} Omega\_0 and omega(W) = -2 Re (W Omega\_0,
c\textasciicircum{}(1)) + O(tau\textasciicircum{}2)), on the one-plaquette
finite graphs H\_FG(G) (|tau| below 48 C\_F) and on the group-G whole-star box
models H\textasciicircum{}G\_N on the open centered boxes Lambda\_N, N at
least 2, each box separately for |tau| below its Kato radius 12
C\_F/(21(2N)\textasciicircum{}3), for G in \{SU(2), SU(3), SU(4), SU(5), U(1),
Z2, SO(3)\}: the Haar moments E[W\textasciicircum{}k], k=1..4, and the
first-order coefficients of omega(W)/tau (1/144, 1/1152, 1/2880, 1/5760, 1/96,
1/48 and 1/864 in that order) agree exactly between the characters route
(forward) and Weyl integration on the SU(N) torus (reverse); the link-flip
lemma holds on these models for SU(2), SU(4), U(1) and Z2, whose central
elements -I, -I, e\textasciicircum{}\{i pi\} and the nontrivial element act as
-1 on the Wilson representation, with E\_3 verified on Lambda\_2, Lambda\_3,
Lambda\_4 and the 24-link coarse factors and E\_2 on [-N,N]\textasciicircum{}2
by exact enumeration; the first-order parity theorem holds for SU(2), SU(4),
SU(5), U(1) and Z2 (E[W\textasciicircum{}3]=0); SU(3) is an obstruction with
E[W\textasciicircum{}3]=1/108, d omega(W\textasciicircum{}2)/dtau=1/6912 and
second-order coefficient 1/589824, SO(3) with 1/27, 1/2592 and 1/331776, and
SU(5) a flip obstruction with fourth-order coefficient
1/63403380965376=1/(2\textasciicircum{}30 3\textasciicircum{}10) on
H\_FG(SU(5)); for SU(2) in the AM2/AQ1 zero-selected patterned family at
kappa=0, |tau| at most 10\textasciicircum{}-8 and both signs,
omega(W\_C)(-tau)=(-1)\textasciicircum{}\{A(C)\} omega(W\_C)(tau) in every
open centered whole-star box and every on-site cutoff, and pointwise for the
limit of the named constructions through the BB2 whole-sequence bound at each
sign (c\textquotesingle{}\_site=2/984375, q=1/64). Under the frozen
convention, the link-flip lemma transfers exactly to the listed gauge groups
that have a central element acting as -1 on the Wilson representation and the
first-order parity theorem exactly to those with vanishing third Haar moment
of W, each with its own first-order Wilson coefficient, on one-plaquette
models and on group box models in each box separately; SU(3) and SO(3) are
recorded obstructions with exact nonzero coefficients, and SU(5) is a recorded
flip obstruction with an exact nonzero fourth-order coefficient; for SU(2) in
the zero-selected family at kappa=0 the Wilson-loop mean is odd or even in tau
according to the parity of the number of plaquettes of any spanning surface of
the loop, in every open centered whole-star box and for the limit of the named
constructions; these are statements about named models and one-plaquette
finite graphs, not AM2 or continuum statements for other groups, not
statements for periodic boxes with an odd side, not uniqueness of any ground
state, and not a statement uniform in the lattice spacing a.
\end{gatequote}

\subsubsection{The frozen convention and the models}
\label{sec:bd1-convention}
Energies are in units $\delta=\alpha/8$. The per-link electric term is
$8C_2$, with $C_F=(N^2-1)/(2N)$ for SU($N$), $C_2(n)=n^2$ for U(1),
$C_2(l)=l(l+1)$ for SO(3) and $C_2=1$ on the odd Z2 state. A face whose four
links carry the Wilson representation $F$ therefore has energy $32C_F$. The
magnetic term is $-(\tau/3)W$ with $W=\operatorname{Re}\chi_F/\dim F$. With the
AM2/AV1 creation sign $\psi=e^{-C}\Omega_0$,
\begin{equation}
 \begin{gathered}
 c^{(1)}=-\frac{\tau/3}{32C_F}\sum_fW_f\Omega_0,\qquad a=\frac{1/3}{32C_F},\\
 \omega(W)=-2\operatorname{Re}\bigl(W\Omega_0,c^{(1)}\bigr)+O(\tau^2)
 =2a\,\mathbb E[W^2]\,\tau+O(\tau^2).
 \end{gathered}
 \label{eq:bd1-first-order}
\end{equation}
The inherited AW1 display $2(W\Omega_0,c^{(1)})$ has the opposite sign; the
contract records it, and both producers reject it (contract defect C7). The
Z2 normalization is the frozen one, $8C_2=8$ on the odd link state; the
readings with $C_2(\mathrm{odd})=1/8$ (coefficient $1/6$) and with
$1-\sigma^x$ (coefficient $1/12$) are rejected controls.

The models are named. The one-plaquette finite graph is
$H_{FG}(G)=32C_2-(\tau/3)W$ on class functions of one plaquette holonomy
(\recid{model_is_finite_graph} true, \recid{transfers_to_aq} false). The
group box model $H^G_N=\sum_e8C_2(e)-(\tau/3)\sum_fW_f$ lives on the open
centered whole-star box $\Lambda_N$, $N\ge2$, with the AQ1 geometry: coarse
24-link factors, the 21 omitted faces of an anchor retained when its whole
star lies in the box, and the selected faces at coefficient 0. For SU(2)
this is exactly the AW1 box Hamiltonian of the zero-selected family. For
every other group it carries no AM2, AV1 or AQ statement.

\paragraph{The Kato radius.} For every listed group $C_F$ is the smallest
nonzero Casimir, attained exactly by $F$ and $\bar F$. The free gap is
therefore $8C_F$ in a box (one excited link) and $32C_F$ on the plaquette,
and $\norm{W_f}\le1$. A simple isolated ground eigenvalue persists,
analytically in complex $\tau$, for
\begin{equation}
 |\tau|<48C_F\quad\text{on }H_{FG}(G),\qquad
 |\tau|<\frac{12C_F}{21(2N)^3}\quad\text{on }H^G_N,
 \label{eq:bd1-kato}
\end{equation}
with $21(2N)^3$ retained faces: $3/448$ for SU(2) at $N=2$. This radius is
sufficient, not sharp, and it shrinks like $N^{-3}$, so nothing is claimed
for all $N$ at once for any group other than SU(2), whose AM2 regime
$|\tau|\le10^{-8}$ is admitted separately. The contract's word ``unique'' in
items~1 and~6 means this simplicity of a finite-box ground eigenvalue, used as
a hypothesis (contract defect C3); it is not uniqueness of any ground state.

\subsubsection{\texorpdfstring{Criterion A: a central element acting as $-1$}{Criterion A: a central element acting as -1}}
\label{sec:bd1-flip}
Let $z$ be central in $G$ with $\rho_F(z)=-I$, and let $E$ be a set of links
meeting every plaquette an odd number of times. Multiplying $U_e$ by $z$ on
every link of $E$ defines a unitary $U_E$. Haar measure is invariant, so
$U_E$ is unitary and fixes $\Omega_0$. Because $z$ is central, $U_E$ commutes
with every translation of each link variable, hence with every endpoint gauge
action; by Schur's lemma it acts on each Peter--Weyl block as a scalar, hence
commutes with every link Casimir and every on-site cutoff projection. A
flipped link multiplies the holonomy of each plaquette through it by $z$ or
$z^{-1}$, which reverses $\chi_F$, and an odd number of flipped links per
plaquette gives $U_EW_fU_E^*=-W_f$. Hence (\projecttag{HNM-BD1-F05})
\begin{equation}
 U_E\,H^G_N(\tau)\,U_E^*=H^G_N(-\tau)\quad\text{in every box and cutoff},\qquad
 U\,H_{FG}(G)(\tau)\,U^*=H_{FG}(G)(-\tau),
 \label{eq:bd1-flip}
\end{equation}
and $\omega_{-\tau}(W)=-\omega_\tau(W)$ wherever the ground eigenvalue is
simple, that is, inside \eqref{eq:bd1-kato}. Such an element exists for
exactly four of the seven groups (Table~\ref{tab:bd1-centre}).

\begin{table}[htbp]
\centering\small
\caption{Central elements acting as $-1$ on the Wilson representation.}
\label{tab:bd1-centre}
\begin{tabular}{llll}
\toprule
Group & Centre & Element with $\rho_F(z)=-I$ & Reason\\
\midrule
SU(2) & $\{I,-I\}$ & $-I$ & $\det(-I)=(-1)^2=1$\\
SU(3) & $\mathbb Z_3$ & none & cube roots of unity are never $-1$\\
SU(4) & $\mathbb Z_4$ & $-I=e^{2\pi i\cdot2/4}I$ & $\det(-I)=(-1)^4=1$\\
SU(5) & $\mathbb Z_5$ & none & fifth roots of unity are never $-1$\\
U(1) & U(1) & $e^{i\pi}$ & charge-1 character value $-1$\\
Z2 & Z2 & the nontrivial element & sign character value $-1$\\
SO(3) & trivial & none & $-I\notin$ SO(3), since $\det(-I)=-1$\\
\bottomrule
\end{tabular}
\end{table}

The absence of a central $-1$ is recorded as the reason, not as the proof,
of non-transfer. The proof is an exact nonzero even-order coefficient
(Section~\ref{sec:bd1-obstructions}): any unitary that commutes with $32C_2$
and reverses $W$ on $H_{FG}(G)$, central or not, would make $\omega(W)$ odd in
$\tau$ and kill every even coefficient.

\subsubsection{Criterion B: the third Haar moment}
\label{sec:bd1-parity}
\paragraph{Single-occurrence orthogonality.} In every listed group a product
of face variables has Haar mean 0 if some link lies in exactly one of the
faces, because as a function of that link the product is a combination of
matrix coefficients of $F$ and $\bar F$. Two distinct plaquettes share at most
one link.

\paragraph{The degenerate level.} The gauge-invariant eigenspace of
$\sum8C_2$ at energy $32C_F$ is spanned by $\chi_F(U_g)\Omega_0$ and
$\overline{\chi_F(U_g)}\Omega_0$ over the plaquettes $g$ of the box. Each
excited link costs at least $C_F$, so at most four links are excited; gauge
invariance forbids a vertex of degree one; $\mathbb Z^3$ is bipartite, so the
excited links form a 4-cycle, which is a plaquette. For U(1) and SU($N$),
$N\ge3$, the level therefore contains the $\operatorname{Im}\chi_F$
directions as well as $W_g\Omega_0$.

\paragraph{The criterion.} On $H_{FG}(G)$ and on every $H^G_N$ inside its
Kato radius, if $\mathbb E[W^3]=0$ then the first-order $\tau$-derivatives at 0
of $\omega(W^2)$, of the centred Euclidean correlation $C(s)$ for every
$s\ge0$ and of the centred real-time correlation $c(\theta)$ for every real
$\theta$ vanish, and the first-order splitting of the whole level vanishes. Every first-order term (state, Duhamel, energy and
vector centring) is a Haar mean of three face variables, which reduces to
$\mathbb E[W^3]$ by single-occurrence orthogonality. Since
$\mathbb E[W^3]=(2\dim F)^{-3}\sum_a\binom3a\mathbb E[\chi^a\bar\chi^{3-a}]$ is a
positive combination of nonnegative integers, $\mathbb E[W^3]=0$ exactly when
every cubic moment vanishes, which also kills the splitting on the complex
directions. Conversely, on $H_{FG}(G)$,
\begin{equation}
 \frac{d\,\omega(W^2)}{d\tau}\Big|_0=2a\,\mathbb E[W^3],\qquad
 C(s)^{(1)}=e^{-32C_Fs}\,\mathbb E[W^3]\bigl(2a+\tfrac s3\bigr),
 \label{eq:bd1-parity}
\end{equation}
which are nonzero for every $s\ge0$ when $\mathbb E[W^3]\ne0$. The vanishing
third moment holds for SU(2), SU(4), SU(5), U(1) and Z2. The flip and the
parity are separate columns: SU(5) has the parity without the flip.

\subsubsection{Moments and first-order coefficients}
\label{sec:bd1-moments}
Table~\ref{tab:bd1-moments} gives the Haar moments and the first-order
coefficients of $\omega(W)/\tau$ from \eqref{eq:bd1-first-order}. For
SU($N$), $N\ge3$, only $\mathbb E[|\chi_F|^2]=1$ survives in $\mathbb E[W^2]$, so
$\mathbb E[W^2]=1/(2N^2)$ and the coefficient is $1/(48N(N^2-1))$. The SU(2)
value $1/144$ reproduces the admitted AW1 coefficient; it is a check of the
convention, not a transferred value, and no SU(2) value enters any other
cell. The first-order coefficients carry the tier
\recid{exact_first_order}; the moments and the higher-order coefficients
carry a route and no tier. The same coefficients hold on the box models,
because every retained face other than $W$ shares at most one link with it.

\begin{table}[htbp]
\centering\small
\caption{Haar moments $\mathbb E[W^k]$, $k=1,\ldots,4$, and first-order Wilson
coefficients under the frozen convention, identical on both routes.}
\label{tab:bd1-moments}
\begin{tabular}{lcccccc}
\toprule
Group & $\mathbb E[W]$ & $\mathbb E[W^2]$ & $\mathbb E[W^3]$ & $\mathbb E[W^4]$ & $C_F$ & $\omega(W)/\tau$ at first order\\
\midrule
SU(2) & 0 & 1/4 & 0 & 1/8 & 3/4 & 1/144\\
SU(3) & 0 & 1/18 & 1/108 & 1/108 & 4/3 & 1/1152\\
SU(4) & 0 & 1/32 & 0 & 7/2048 & 15/8 & 1/2880\\
SU(5) & 0 & 1/50 & 0 & 3/2500 & 12/5 & 1/5760\\
U(1) & 0 & 1/2 & 0 & 3/8 & 1 & 1/96\\
Z2 & 0 & 1 & 0 & 1 & 1 & 1/48\\
SO(3) & 0 & 1/9 & 1/27 & 1/27 & 2 & 1/864\\
\bottomrule
\end{tabular}
\end{table}

The cubic moments $\mathbb E[\chi^a\bar\chi^b]$, $a+b=3$, all vanish except for
SU(3), where the invariant $\varepsilon$ tensor in $F^{\otimes3}$ gives
$\mathbb E[\chi^3]=\mathbb E[\bar\chi^3]=1$, and SO(3), where the real vector
representation carries $\varepsilon_{ijk}$ and every cubic moment is 1. The
forward also exports $\mathbb E[W^5]$ ($5/1296$ for SU(3), $1/50000$ for SU(5),
$2/81$ for SO(3), zero otherwise).

\subsubsection{The obstruction cells}
\label{sec:bd1-obstructions}
On $H_{FG}(G)$ the second-order coefficient of $\omega(W)$ is
$(\psi_1,W\psi_1)+2(W\Omega_0,\psi_2)=3a^2\,\mathbb E[W^3]$, so a nonzero third
moment gives a nonzero even coefficient.
\begin{itemize}
\item \textbf{SU(3)} ($a=1/128$): $\mathbb E[W^3]=1/108$,
 $d\omega(W^2)/d\tau=1/6912$, and second-order coefficient $1/589824$. The
 level $\{\chi_F,\chi_{\bar F}\}$ splits at first order by $\pm1/18$.
\item \textbf{SO(3)} ($a=1/192$): $\mathbb E[W^3]=1/27$,
 $d\omega(W^2)/d\tau=1/2592$, and second-order coefficient $1/331776$.
\item \textbf{SU(5)} ($a=5/1152$): $\mathbb E[W^3]=0$, so the parity column is a
 transfer. Each application of $W$ changes the 5-ality by $\pm1$, so the
 order-$m$ coefficient involves Haar means of $m+1$ factors $\chi_F$ or
 $\bar\chi_F$ with $a-b$ divisible by 5; orders 1 to 3 give no even
 contribution, and the first even order that can be nonzero is 4. There
 \begin{equation}
 c_4=\frac1{63403380965376}=\frac1{2^{30}3^{10}}\approx1.57720295790\times10^{-14},
 \label{eq:bd1-su5}
 \end{equation}
 by three evaluations: the direct Rayleigh--Schr\"odinger series,
 Hellmann--Feynman $c_4=-15E_5$ with $E_5=-1/951050714480640$, and a closed
 form. The forward's closed form uses the single channel
 $\Lambda^2\leftrightarrow\Lambda^3$ inside $\psi_2$; the reverse's uses the
 column chain $\Lambda^1\to\Lambda^2\to\Lambda^3\to\Lambda^4\to$ trivial, which
 gives $c_4=(10/81)\,10^{-5}/\prod_kE(\Lambda^k)$ with
 $E(\Lambda^k)=384/5,\,576/5,\,576/5,\,384/5$. SU(5) is therefore a flip
 obstruction with the parity column a transfer.
\end{itemize}
The one-plaquette series through order 5 are exported for all seven groups;
the SU(2) row reproduces the admitted AW1 fixture
$\tau/144-(5/11943936)\tau^3$. As a supplementary argument, not used in place
of the coefficients, a unitary reversing $W$ on $H_{FG}(G)$ needs the range of
$W$ to be symmetric; it is $[-1/2,1]$ for SU(3) and $[-1/3,1]$ for SO(3), and
$-1$ is not attained for SU(5).

\subsubsection{The transfer ledger}
\label{sec:bd1-ledger}
Table~\ref{tab:bd1-ledger} is the ledger both producers export. Every
transfer is to $H_{FG}(G)$ and $H^G_N$: the operator identity in every box and
cutoff, and the oddness and first-order statements in each box inside its
Kato radius.

\begin{table}[htbp]
\centering\small
\caption{The BD1 transfer ledger. T: transfer to the named models with its own
exact checks. O: obstruction, with its exact counterexample on
$H_{FG}(G)$.}
\label{tab:bd1-ledger}
\begin{tabularx}{\linewidth}{@{}l>{\raggedright\arraybackslash}X>{\raggedright\arraybackslash}X ll@{}}
\toprule
Group & Flip lemma & Parity theorem & First order & AM2/AQ chain, dictionary\\
\midrule
SU(2) & T ($-I$; the admitted AW1 lemma) & T & 1/144 (check) & admitted reference\\
SU(3) & O: $c_2=1/589824$ (reason $\mathbb Z_3$) & O: $d\omega(W^2)/d\tau=1/6912$ & 1/1152 & obligation\\
SU(4) & T ($-I$) & T & 1/2880 & obligation\\
SU(5) & O: $c_4=1/(2^{30}3^{10})$ (reason $\mathbb Z_5$) & T & 1/5760 & obligation\\
U(1) & T ($e^{i\pi}$) & T & 1/96 & obligation\\
Z2 & T (nontrivial element) & T & 1/48 & obligation\\
SO(3) & O: $c_2=1/331776$ (trivial centre) & O: $1/2592$ & 1/864 & obligation\\
\bottomrule
\end{tabularx}
\end{table}

For every group other than SU(2) the AM2/AQ chain is an obligation: the AM2
re-instantiation with its finite-volume ground state and gap, the AV1
product split, the AQ1 construction and dynamics with the AQ2 gap, the BB1
and BB2 comparisons, and a dictionary to a bare coupling. The review records these as rows
O1--O7 and O9 of its obligations table; row O8 covers F2 and literal vertex
boxes, O10 SU(2) at a nonzero selected triple and O11 periodic boxes with an
odd side.

\subsubsection{Flip sets}
\label{sec:bd1-flipsets}
The flip set is
$E_3=\{(p,x):p_y\text{ even}\}\cup\{(p,y):p_z\text{ even}\}\cup\{(p,z):p_x\text{ even}\}$,
the admitted AW1 set, and in two dimensions $E_2=\{(p,x):p_y\text{ even}\}$.
The reverse derives $E_3$ before reading the contract: of the 27 rules ``link
$(p,d)$ in $E$ if and only if $p_{\sigma(d)}$ is even'', exactly two are odd on
all 24 plaquette classes, the cyclic rule, which is $E_3$, and the
anticyclic rule, and they differ by a centre gauge transformation.
Table~\ref{tab:bd1-flipsets} gives the enumeration on the boxes.

\begin{table}[htbp]
\centering\small
\caption{$E_3$ on the open centered whole-star boxes: every owned plaquette
and every retained face meets $E_3$ an odd number of times.}
\label{tab:bd1-flipsets}
\begin{tabular}{rrrrrr}
\toprule
$N$ & Owned links & Plaquettes & Met once & Met three times & Retained faces (once / three times)\\
\midrule
2 & 3000 & 2335 & 1082 & 1253 & 1344 (672 / 672)\\
3 & 8232 & 6909 & 3630 & 3279 & 4536 (2268 / 2268)\\
4 & 17496 & 15291 & 7348 & 7943 & 10752 (5376 / 5376)\\
\bottomrule
\end{tabular}
\end{table}

The $N=2$ counts reproduce the AW1 forward report. On seven probed coarse
factors every one of the 52 faces meeting the factor (49 omitted, 3
selected) meets $E_3$ oddly; $E_3$ holds 16 links of a factor with even $z$
and 8 with odd $z$, so it is not invariant under odd coarse translations in
$z$, which the flip lemma does not need. $E_2$ meets every plaquette of
$[-N,N]^2$ exactly once (16, 36 and 64 plaquettes for $N=2,3,4$).

\paragraph{Periodic tori: the GF(2) reading.} On periodic tori with an odd
side, $E_3$ and $E_2$ meet seam plaquettes evenly: 16 even plaquettes on
$3\times4\times4$, $4\times3\times4$ and $4\times4\times3$, 24 on
$3\times3\times4$, 27 on $3\times3\times3$, and 4 on the $4\times3$ torus. $E_2$,
keyed to $p_y$ only, still meets every plaquette of the $3\times4$ torus once.
Gaussian elimination over GF(2) shows more: a link set odd on every plaquette
of a torus exists exactly when every coordinate plane has an even plaquette
count, that is, when at most one side is odd in three dimensions, or when
$L_xL_y$ is even in two. The contract's periodic ``obstruction'' therefore
concerns $E_3$ and $E_2$ at the seam, not every flip set (contract defect
C4). These are recorded findings, not claims: no flip or parity statement is
made on any periodic box with an odd side.

\subsubsection{SU(2) area parity}
\label{sec:bd1-area}
For SU(2) in the zero-selected family at $\kappa=0$, $|\tau|\le10^{-8}$ and
both signs, let $W_C$ be the half-trace of the holonomy of a closed path $C$
of plaquette edges and $A(C)$ the number of plaquettes of a spanning
surface. Every plaquette meets $E_3$ oddly and every interior link of a
surface is counted twice, so for every spanning surface $S$ of $C$,
$|S|\equiv|C\cap E_3|\pmod2$, independently of $S$; in particular every closed
surface of plaquettes has an even number of them. In every open centered
whole-star box, every on-site cutoff and for the untruncated ground vectors,
the AW1 identity $U_EH_N(\tau)U_E^*=H_N(-\tau)$ and AM2 simplicity give
\begin{equation}
 \omega_{N,-\tau}(W_C)=(-1)^{A(C)}\,\omega_{N,\tau}(W_C),
 \label{eq:bd1-area}
\end{equation}
because $-1$ is central and each traversal of a flipped link contributes one
factor $-1$. For the limit of the named constructions the identity holds
pointwise at each sign: $U_E$ is a product over links, so
$\rho^{F1,N}_Y(-\tau)=U_{E\cap Y}\rho^{F1,N}_Y(\tau)U_{E\cap Y}^*$ for the region
$Y$ owning the links of $C$, and the BB2 whole-sequence bound
\eqref{eq:bb2-cauchy} with $c'_{\rm site}=2/984375$ and $q=1/64$ carries it to
the limit at each sign separately. More generally
$\omega^{-\tau}_\infty=\omega^{\tau}_\infty\circ\alpha_E$. The AW1 whole-set
statement alone would give this only as sets of subsequential limits, and
using it for a pointwise claim is a rejected mutation. Observables fixed by
$U_E$, such as $W_C^2$, are even in $\tau$.

The surface identity is audited on nine loops with two disjoint spanning
surfaces each and five closed surfaces (6, 10, 14, 24 and 32 plaquettes,
including a ring of eight cubes of Euler characteristic 0) in the forward, on
6480 rectangles and 400 random surfaces in the reverse, and on 216 planar
loops in the skeptic's pre-comparison check. The audits check the identity;
they do not replace the proof. The statement is made only for SU(2), the
zero-selected family, $\kappa=0$, open centered whole-star boxes, their
cutoffs and the limit of the named constructions; not for periodic boxes with
an odd side, F2 or literal vertex boxes, nonzero selected triples or other
groups.

\subsubsection{Two routes and the review}
\label{sec:bd1-review}
The forward route computes each moment three ways inside the character route
(iterated action on the character basis, iterated tensor multiplicities, and
closed-form counts by Schur--Weyl duality with the hook-length formula,
charge balance, parity and Catalan numbers) and adds a labelled internal Weyl
cross-check that is never an admission value. The reverse route never uses a
character table: it integrates on the SU($N$) torus by constant terms
against the Weyl density, with exponents taken modulo the diagonal, on the
SO(3) torus with density $(2-z-z^{-1})/2$, by the binomial constant term for
U(1) and by summation for Z2, and it applies the electric operator to Weyl
numerators. The forward checker runs 37 exact checks and the reverse 64,
each with all 22 controls as damaging mutations (111 and 110 rejections);
both replay byte for byte under \texttt{-B} and \texttt{-B -O}.

The skeptic computed every value by three routes of its own (characters,
Weyl constant terms, and Frobenius and hook-length counts). The moments for
$k=1,\ldots,4$, the first-order coefficients and derivatives, the
one-plaquette series and the SU(5) value $\mathbb E[W^5]$ agree exactly across
forward, reverse and skeptic, 141 cells in all, with no single-route cell.
The review records a trap it tested. On the U($N$) torus instead of the
SU($N$) torus the SU(3) third moment would be 0 and the SU(4) fourth moment
$3/1024$; both producers reduce modulo the diagonal, return $1/108$ and
$7/2048$, and abort under a U($N$) substitution. The reverse committed at
05:53:54 UTC and the forward at 05:54:17. The skeptic's values were final at
05:45:05, before either producer commit, and its package was committed last,
at 06:10:08; its exposure before that was to file names, hashes and possibly
commit subject lines. The post-comparison checker runs 29 checks and 71
source-edit runs: 44 validator weakenings, 17 must-abort substitutions, and 6
silent edits, of which 2 are harmless and 4 are caught only by the skeptic's
validator. It also rejects 27 damaged packets. Both closures of 27 files
verify, each with its 23 declared inputs. The review found no blocking issue.
It records a wrong preview digit string for $c_4$ in the reverse report,
whose results carry the exact value; the forward's literal flip-scope field
beside its named groups; two reverse fields outside the contract list; and
both reports written through the shell after a tool refusal.

\paragraph{Contract defects.} Seven are recorded, none blocking: the
preregistered tier rule conflicts with the tier control for second- and
fourth-order coefficients, and the control governs (C1); the frozen flip
scope field is a description, so the gate names the groups (C2); ``unique''
means finite-box simplicity (C3); the periodic obstruction concerns $E_3$ and
$E_2$ at the seam (C4); the two-route table comparison is a pair-level
comparison in a paired loop (C5); the contract does not state that the box
coefficient equals the plaquette value, which both producers prove (C6); and
the inherited AW1 display sign (C7).

\subsubsection{Limitations recorded by the gate}
\label{sec:bd1-limits}
The BD1 gate records the following limitations verbatim.
\begin{gatequote}
\begin{itemize}
\item Scope: statements about named finite models only, H\_FG(G) (finite
graphs; model\_is\_finite\_graph true, transfers\_to\_aq false) and
H\textasciicircum{}G\_N box by box inside a Kato radius that shrinks like
N\textasciicircum{}-3; nothing is claimed for all N at once for any group
other than SU(2), no limit state is constructed for any other group, and no
AM2, AV1, AQ, BB1 or BB2 statement and no dictionary to a bare coupling is
transferred to SU(3), SU(4), SU(5), U(1), Z2 or SO(3); each is recorded as an
obligation.
\item The group-cell oddness statements use simplicity of the finite-box
ground eigenvalue from Kato\textquotesingle{}s theorem with a sufficient, not
sharp, radius (12 C\_F/F\_N with F\_N=21(2N)\textasciicircum{}3 retained
faces; 48 C\_F on one plaquette); this is simplicity in a finite box, not
uniqueness of any ground state, and the word "unique" in contract items 1 and
6 is read that way.
\item No flip or parity statement is made on any periodic box with an odd
side. The recorded obstruction is that E\_3 and E\_2 meet seam plaquettes
evenly when the keyed coordinate has an odd side (16 on 3x4x4, 4x3x4 and
4x4x3; 24 on 3x3x4; 4 on the 4x3 torus); E\_2 still meets every plaquette of
the 3x4 torus once, and by GF(2) elimination some flip set exists whenever
every coordinate plane has an even plaquette count (at most one odd side);
these are recorded findings, not claims.
\item The SU(2) corollary inherits the AW1 flip lemma, AM2 simplicity and
cutoff removal and the BB2 whole-sequence convergence with their own
limitations (the named constructions F1 and F2 only; no statement about other
states or boundary conditions); it is made at kappa=0 only, not at a nonzero
selected triple and not for F2 or literal vertex boxes; the loop and surface
enumerations audit the identity |C cap E\_3| = A(C) mod 2 and do not replace
the proof.
\item Exact rational arithmetic decides every value; decimals are truncated
previews (the reverse report prints 1.57720727258e-14 for omega\_4 of SU(5), a
text slip; the exact value 1/63403380965376 and the results preview
1.57720295790e-14 stand). Kato\textquotesingle{}s theorem, Peter-Weyl, the
Weyl integration and character formulas, Schur-Weyl duality and
Schur\textquotesingle{}s lemma are cited, not machine-checked; the finite
enumerations and fixtures audit their finite consequences.
\item Contract defects recorded, none blocking: the preregistered
tier\_label\_rule conflicts with the tier\_mixing\_rejected control for
second- and fourth-order coefficients (both producers follow the control:
route named, no tier); gate\_fields\_required flip\_transfer\_scope is a
description, so the gate names SU(2), SU(4), U(1) and Z2; items 1 and 6 say
"unique" for finite-box simplicity; the periodic-torus obstruction concerns
E\_3 and E\_2 at the seam, not every flip set.
\item Review standing: model-agent skeptic with correlated ancestry (same
model family as both producers and the advisor); its pre-comparison package
was committed after both producer packets, with its values final before either
producer commit; not human peer review; scientific priority is not verified;
weak coupling and the continuum are not addressed, and no estimate is made
uniformly in the lattice spacing a.
\end{itemize}
\end{gatequote}
Full derivations and reviewed implementations:
\repo{research/round33/forward/bd1/report.md}{BD1 forward},
\repo{research/round33/reverse/bd1/report.md}{BD1 reverse},
\repo{research/round33/skeptic/bd1.md}{independent model-agent review} and the
\repo{research/round33/contracts/bd1.json}{frozen contract}; the transferred
statements are those of the \repo{research/round32/advisor/aw1-gate.json}{AW1
gate}, which is not edited.

\subsection{BD2: new observables and a new dimension in SU(2)}
\label{sec:bd2}

\textbf{Contribution \projecttag{HNM-C-BD2}.} The BD2 contract, recorded
under the title ``Applications II: new observables and a new dimension in
SU(2) - the 1x2 Wilson loop, the electric energy band and the two-plaquette
graph with independent couplings, and the 2+1D AM2 constants'', applies
admitted SU(2) equations to four new cases
(\repo{research/round33/advisor/selection-bd2.md}{selection note}). It is
investigation~8 of~8 and single-direction: one forward producer, with 33
shared premises, and the skeptic's replay from the contract. Its blind
discriminating target is the relative width of the certified enclosure of
the 1x2 loop on the graph at $D=8$, at most $1/10^{16}$; the other items are
exact computations and format checks. It lists 21 controls. Its gate is
recorded under the title ``Hruday new observables and dimensions in SU(2):
the two-plaquette graph, the 1x2 loop, the electric band and 2+1D''.

\subsubsection{Reviewed scope}
\label{sec:bd2-scope}
The \repo{research/round33/advisor/bd2-gate.json}{BD2 gate} records the
verdict \recid{accepted_within_scope}, with sub-label
\recid{sign_certified_finite_graph} and secondary sub-labels
\recid{transfer_to_named_model}, \recid{obstruction_recorded} and
\recid{static_not_dynamic}. Its gate fields set
\recid{graph_sign_certified}, \recid{z3_1x2_formal_only},
\recid{electric_band_claimed} and \recid{model_is_finite_graph} to true, with
the band scope ``omega(h\_R) for every finite box of F1 and F2 (untruncated at
fixed N) and the limit of the named constructions, both signs''. Seven
fields are false: \recid{area_law_claimed}, \recid{transfers_to_aq},
\recid{uniqueness_of_ground_state_claimed}, \recid{continuum_claim},
\recid{rate_in_a_claimed}, \recid{weak_coupling_claim} and
\recid{scientific_priority_verified}. The accepted statement, quoted
verbatim below in the gate's ASCII notation, is the scope of every claim in
this subsection.
\begin{gatequote}
On the Round11 two-plaquette graph H\_FG(l1,l2) = K - l1 W\_1 - l2 W\_2 (open
two-square patch, 6 vertices, 7 links, gauge-invariant sector, K the sum of
the seven link Casimirs at rho=1, alpha units, trace-monomial cutoffs D=6 and
D=8), the coefficients through total order 4 are identical at both cutoffs:
<W\_1> = l1/6 - 5 l1\textasciicircum{}3/864 + l1 l2\textasciicircum{}2/4212,
<z> = 7 l1 l2/216 - 349 (l1\textasciicircum{}3 l2 + l1
l2\textasciicircum{}3)/303264 for the 1x2 loop z, and <C\_shared> =
(l1\textasciicircum{}2 + l2\textasciicircum{}2)/48 - 5 (l1\textasciicircum{}4
+ l2\textasciicircum{}4)/4608 - 49 l1\textasciicircum{}2
l2\textasciicircum{}2/438048; <W\_1> is odd in l1 and even in l2 by the
one-link flips of h1 and h2, <z> is odd in each coupling and <C\_shared> even
in each; at l1=l2=l with |l| in \{1/10, 1/100, 1/1000\} and both signs the
certified enclosures of <z> are positive at D=6 and D=8, with D=8 relative
width at most 1.79645e-17 against the frozen 1/10\textasciicircum{}16. In the
zero-selected SU(2) family (AM2/AQ1, |tau| at most 10\textasciicircum{}-8,
cover R=\{0,e\_z\}) the 1x2 Wilson loop W\_\{1x2\} of the omitted xz faces at
0 and e\_x has first-order coefficient 0 in every box, satisfies
|omega(W\_\{1x2\})| <= K\_2\textquotesingle{} tau\textasciicircum{}2
(K\_2\textquotesingle{} the AY1 gate constant; about 1.34175e-12 at
|tau|=10\textasciicircum{}-8) for every F1 and F2 box and for the limit of the
named constructions at both signs, is even in tau in every F1 box and for the
limit, and has the formal second-order coefficient 7/124416 (labelled
formal\_second\_order\_coefficient: no sign and no certified remainder); the
electric energy omega(h\_R) on R lies in [(3/2)L\textasciicircum{}2, 98|tau|],
about [2.87586637818e-19, 9.8e-7] at |tau|=10\textasciicircum{}-8, with L the
AY2 lower distance, for every F1 and F2 box (untruncated at fixed N) and for
the limit, at both signs. In the declared 2+1-dimensional model the AM2
counting gives p=3, J=|tau|, termination order 6, G\_3(R)=9
e\textasciicircum{}\{3/32\} and G\_3\textquotesingle{}(R)=118
e\textasciicircum{}\{3/32\} at R=1/64, and the cap |tau| <= 1/(576
e\textasciicircum{}\{3/32\}) (directed lower bound
1580747155173/10\textasciicircum{}15), with no finite-volume theorem. On the
Round11 two-plaquette graph with independent couplings, the exact low-order
coefficients of the face mean, the 1x2 loop and the shared-link Casimir are
computed with zero truncation error, the face mean is odd in its own coupling
and even in the other, and the 1x2 loop has a certified sign at every grid
point; in the zero-selected SU(2) family the 1x2 Wilson mean vanishes at first
order and is even in tau, with only a formal second-order coefficient, the
electric energy on the cover lies in a certified two-sided band for every box
and for the limit of the named constructions, and the 2+1-dimensional AM2
constants are recounted; these are statements about named models and a finite
graph, not an area law, not a certified sign of the 1x2 mean, not uniqueness
of any ground state, and not a statement uniform in the lattice spacing a.
\end{gatequote}

\subsubsection{The two-plaquette graph with independent couplings}
\label{sec:bd2-graph}
The graph is the Round11 open two-square patch \cite{r11}: vertices TL, TM, TR, BL, BM,
BR and seven links $h_1$ (TL$\to$TM), $h_2$ (TM$\to$TR), $h_3$ (BL$\to$BM),
$h_4$ (BM$\to$BR), $v_L$, $v_M$ and $v_R$, with holonomies
$U=v_Mh_3^{-1}v_L^{-1}h_1$ and $V=h_2v_Rh_4^{-1}v_M^{-1}$, face variables
$x=W_1=\tfrac12\operatorname{Tr}U$ and $y=W_2=\tfrac12\operatorname{Tr}V$, and the 1x2
loop $z=\tfrac12\operatorname{Tr}(UV)$, in which the shared link $v_M$ cancels.
In alpha units, on the gauge-invariant sector,
\begin{equation}
 H_{FG}(l_1,l_2)=K-l_1W_1-l_2W_2,\qquad K=3C_U+3C_V+C_{\rm shared},
 \label{eq:bd2-graph}
\end{equation}
with $K$ the sum of the seven link Casimirs at $\rho=1$ and $C_{\rm shared}$
the Casimir of $v_M$. AZ2 treated the equal couplings $l_1=l_2$
(\repo{papers/round32-addendum/main.pdf}{Round32 addendum}); here the two
faces are coupled independently.

\paragraph{Exact coefficients.} $K$ is triangular in the degree, so every
Rayleigh--Schr\"odinger vector $\psi_{mn}$ lies in the polynomials of degree
$m+n$ and the equations hold as polynomial identities for $m+n\le5$. The
truncation error is therefore exactly 0, and the tables computed at the
trace-monomial cutoffs $D=6$ and $D=8$ are identical. The first vectors are
$\psi_{10}=x/3$, $\psi_{01}=y/3$ and $\psi_{11}=4xy/39+4z/351$. Through total
order 4,
\begin{equation}
 \begin{aligned}
 \langle W_1\rangle&=\frac{l_1}6-\frac{5l_1^3}{864}+\frac{l_1l_2^2}{4212}+\cdots,\\
 \langle z\rangle&=\frac{7l_1l_2}{216}-\frac{349(l_1^3l_2+l_1l_2^3)}{303264}+\cdots,\\
 \langle C_{\rm shared}\rangle&=\frac{l_1^2+l_2^2}{48}-\frac{5(l_1^4+l_2^4)}{4608}-\frac{49l_1^2l_2^2}{438048}+\cdots,\\
 E_0&=-\frac{l_1^2+l_2^2}{12}+\frac{5(l_1^4+l_2^4)}{3456}-\frac{l_1^2l_2^2}{8424}+\cdots,
 \end{aligned}
 \label{eq:bd2-series}
\end{equation}
with every other coefficient of total order at most 4 exactly zero. The
checks are exact: Hellmann--Feynman
$\langle W_1\rangle_{(m,n)}=-(m+1)E_{(m+1,n)}$ through total order 5; a second
route for $\langle C_{\rm shared}\rangle$ at second order, from
$E_{20}(\rho)=-1/(3(3+\rho))$ and its $\rho$-derivative $1/48$; the section
$l_2=0$, which gives the one-face value $-5/864$; and the diagonal $l_1=l_2$,
which reproduces the admitted AZ2 series ($1/6$ and $-187/33696$ for
$\langle W_1\rangle$, $7/216$ and $-349/151632$ for $\langle z\rangle$). At
$l_1=l_2$ the second-order coefficients are $7/216$ for $\langle z\rangle$ and
$1/24$ for $\langle C_{\rm shared}\rangle$.

\paragraph{Flip parities.} For a link $e$ the flip $U_e\mapsto-U_e$ commutes
with every Casimir, with the six vertex gauge actions and with Haar measure,
and reverses every loop through $e$. Table~\ref{tab:bd2-flips} lists the
actions. The ground state of $H_{FG}$ is simple for every real
$(l_1,l_2)$ (Round11 Theorem~F), so the flip of $h_1$, a non-shared link of
face~1, makes $\langle W_1\rangle$ and $\langle z\rangle$ odd in $l_1$ and
$\langle C_{\rm shared}\rangle$ even, and the flip of $h_2$, a non-shared link
of face~2, makes $\langle W_1\rangle$ even in $l_2$, $\langle z\rangle$ odd and
$\langle C_{\rm shared}\rangle$ even. Coefficient by coefficient,
$F_{h_1}\psi_{mn}=(-1)^m\psi_{mn}$ and $F_{h_2}\psi_{mn}=(-1)^n\psi_{mn}$. The
shared-link flip alone gives only the joint statement, and reading a
two-variable parity from it, or checking only at $l_1=l_2$, is a rejected
control.

\begin{table}[htbp]
\centering\small
\caption{One-link centre flips on the two-plaquette graph.}
\label{tab:bd2-flips}
\begin{tabular}{lll}
\toprule
Flipped link & Action on $(x,y,z)$ & Effect on $H_{FG}(l_1,l_2)$\\
\midrule
$h_1$, $v_L$ or $h_3$ (face 1, non-shared) & $(-x,y,-z)$ & $H(-l_1,l_2)$\\
$h_2$, $v_R$ or $h_4$ (face 2, non-shared) & $(x,-y,-z)$ & $H(l_1,-l_2)$\\
$v_M$ (shared) & $(-x,-y,z)$ & $H(-l_1,-l_2)$ only\\
\bottomrule
\end{tabular}
\end{table}

\paragraph{Certified enclosures of $\langle z\rangle$.} At $l_1=l_2=l$ the
Hamiltonian is the AZ2 model, and the ground-state certificate is the
admitted AZ2 certificate code, copied verbatim (63 blocks compared textually
with the snapshot); only the observable changes. A rational Ritz vector with
complete residual, the Weyl gap $E_1-E_0\ge3-2|l|$, Temple's bound, the Round11
Schur-complement tail comparison, and the Eckart and Davis--Kahan angles give
$\sin\theta\le s$, and then
$|\langle v,zv\rangle/\langle v,v\rangle-\langle z\rangle|\le2s\sigma_z+2s^2\norm z$
with $\sigma_z\le\norm z=1$. Ends are rounded outward to denominators
$10^{45}$. Table~\ref{tab:bd2-enclosures} gives the results.

\begin{table}[htbp]
\centering\small
\caption{Certified enclosures of $\langle z\rangle$ at $l_1=l_2$ (both signs
give the same enclosure, since $\langle z\rangle$ is even under the shared-link
flip). Relative width is $(\text{upper}-\text{lower})/\min(|\text{lower}|,|\text{upper}|)$.}
\label{tab:bd2-enclosures}
\begin{tabular}{rllll}
\toprule
$D$ & $|l|$ & Lower end (preview) & Relative width & Sign\\
\midrule
6 & 1/1000 & $3.240740510578259506053328\times10^{-8}$ & $6.79043\times10^{-22}$ & $+$\\
6 & 1/100 & $3.240717724665334529490919\times10^{-6}$ & $6.81088\times10^{-17}$ & $+$\\
6 & 1/10 & $3.238440859055494033169085\times10^{-4}$ & $7.02377\times10^{-12}$ & $+$\\
8 & 1/1000 & $3.240740510578259506054428\times10^{-8}$ & $1.73689\times10^{-31}$ & $+$\\
8 & 1/100 & $3.240717724665334639851747\times10^{-6}$ & $1.74212\times10^{-24}$ & $+$\\
8 & 1/10 & $3.238440859066867050597032\times10^{-4}$ & $1.79645\times10^{-17}$ & $+$\\
\bottomrule
\end{tabular}
\end{table}

Every enclosure excludes 0 with a positive sign, and every $D=8$ enclosure
lies inside its $D=6$ one, which is observed and not a target. The worst
$D=8$ relative width is $1.79645\times10^{-17}$, at $|l|=1/10$, against the
frozen $1/10^{16}$: a margin of about $5.57$. At $l=+1/10$ and $D=8$ the exact
ends are
\begin{equation}
 \begin{gathered}
 \text{lower}=\frac{40480510738335838132462907353125470123319}{125\cdot10^{42}},\\
 \text{upper}=\frac{64768817181337342175480593943907904225707}{2\cdot10^{44}}.
 \end{gathered}
 \label{eq:bd2-ends}
\end{equation}
Each point carries the five-part AZ2-type residual ledger: seven per-link rows
(thresholds $123/2$ and 96 on the non-shared links, 45 and 69 on $v_M$), the
joint product channel with its corner overlaps and the positive interference,
the gauge projection, which is exactly 0 with its reason, the Ritz residual
with both angles and the Weyl gap, and the directed arithmetic. The skeptic
re-evaluated the lemma in its own spin-network basis with its own Ritz
vectors and reproduced every angle and width to 8 digits. Its own sharper
tail certificate gives a worst width of $3.56\times10^{-18}$ (margin
$28.09$), with enclosures inside the producer's, and Davis--Kahan alone would
give $8.72\times10^{-17}$; any of the three meets the target.

\subsubsection{\texorpdfstring{The $\mathbb Z^3$ 1x2 rectangle}{The Z3 1x2 rectangle}}
\label{sec:bd2-rectangle}
In the zero-selected family the frozen rectangle is the union of the omitted
$xz$ faces at 0 and $e_x$, with links $(0,x)$, $(e_x,x)$, $(2e_x,z)$,
$(e_x+e_z,x)$, $(e_z,x)$ and $(0,z)$; the shared middle link cancels. All six
links are owned by $R=\{0,e_z\}$, so $W_{1\times2}\in B(\mathcal H_R)$ with
$\norm{W_{1\times2}}\le1$, and on the seven links of the two faces
$(W_{F1},W_{F2},W_{1\times2})$ is the graph's $(x,y,z)$.

\paragraph{The first-order zero.} Every link of $W_{1\times2}$ occurs once and
a single face covers at most three of its six links, so
$\mathbb E[W_{1\times2}W_f]=0$ for every face and the first-order coefficient is
exactly 0 in every box of F1 and F2.

\paragraph{The bound.} Writing
$\omega(W_{1\times2})=\tr((\rho_R-P_R-\rho^{(1)}_R)W_{1\times2})$, since the
product and the first-order density contribute 0,
\begin{equation}
 |\omega(W_{1\times2})|\le K_2'\tau^2\approx1.34175\times10^{-12}\quad(|\tau|=10^{-8}),
 \qquad K_2'\approx13417.5275439,
 \label{eq:bd2-rectangle}
\end{equation}
at both signs, with $K_2'$ the AY1 gate constant. It holds for every F1 and F2
box through the AY1 per-box ball (AY1 forward F11--F14, a reviewed step of the
admitted proof, passed to the untruncated vector by AV1 F22) and for the limit
of the named constructions through the AY1 and AY2 gate ball and BB2
items~2--3. The bound carries no sign.

\paragraph{Evenness, for F1 and the limit only.} The flip set $E_3$ meets each
of the two faces in 3 links and the rectangle in all 6, so
$U_EW_{1\times2}U_E^*=W_{1\times2}$ (area 2). With the AW1 identity at
$\kappa=0$, $\omega_{N,-\tau}(W_{1\times2})=\omega_{N,\tau}(W_{1\times2})$ in every
open centered whole-star box (F1) and every cutoff, and the limit is even by
BB2 at each sign. Evenness is not claimed for F2 boxes, and it gives no
remainder bound.

\paragraph{The formal second-order coefficient.} Only the ordered pairs of
the two faces cover the rectangle evenly, with
$\mathbb E[W_{F1}W_{F2}W_{1\times2}]=1/16$, and $W_{1\times2}\Omega_0$ has energy
36, so formally
\begin{equation}
 \omega(W_{1\times2})=\Bigl(\frac1{31104}+\frac1{41472}\Bigr)\tau^2+\cdots=\frac7{124416}\,\tau^2+\cdots .
 \label{eq:bd2-formal}
\end{equation}
This is labelled \recid{formal_second_order_coefficient}: it has no sign, no
certified value and no remainder. The identity $7/124416=(7/216)/24^2$
relates it to the graph coefficient of $\langle z\rangle$ under
$\tau_{FG}=\tau/24$; it is a labelled identity, not a transfer. The certified
bound exceeds the formal term by a factor of about $2.38479\times10^8$, so a
certified sign of the $\mathbb Z^3$ 1x2 mean is a recorded obstruction.

\subsubsection{The electric band}
\label{sec:bd2-band}
The observable $h_R=8\sum C_e$ over the 48 links of $R$ is nonnegative and
unbounded; $\omega(h_R)$ is the monotone limit over its spectral cutoffs. For
every F1 and F2 box (untruncated at fixed $N$) and for the limit, at both
signs,
\begin{equation}
 \tfrac32L^2\le\omega(h_R)\le98|\tau|,\qquad
 \tfrac32L^2\approx2.87586637818\times10^{-19},\quad98|\tau|=\frac{49}{50000000}\quad(|\tau|=10^{-8}).
 \label{eq:bd2-band}
\end{equation}

\paragraph{Lower end.} The spectrum of $h_R$ is $\{0\}\cup[6,\infty)$ with
kernel $P_R$, since the smallest nonzero value of $8\sum j_e(j_e+1)$ is
$8\cdot\tfrac34=6$; hence $\omega(h_R)\ge6\tr(Q_R\rho_R)$ with $Q_R=I-P_R$. For a
pure reference the Fuchs--van de Graaf inequality, proved inline, gives
$\tfrac12\norm{\rho_R-P_R}_1\le\sqrt{\tr(Q_R\rho_R)}$, so
$\omega(h_R)\ge6(L/2)^2=\tfrac32L^2$ whenever $\norm{\rho_R-P_R}_1\ge L$. The
distance $L\approx4.37863\times10^{-10}$ is the AY2 gate lower end for the
limit, and the same number from the AY1 per-box ball for every box (tier
\recid{first_order_distance_from_product}).

\paragraph{Upper end.} The AQ1 reset on $F=R$ charges the seven incident
anchors, each whole star of norm $7|\tau|$, giving
$\omega_N(h_R)\le2\cdot7\cdot7|\tau|=98|\tau|$ for every F1 box; for F2 boxes the
same budget is AY1 gate item~(1) (tier \recid{crude_majorant}). The bound passes
to the limit only by lower semicontinuity through the cutoff, and an exact
fixture shows why the other direction fails: $\rho_k=(1-1/k)P_0+(1/k)P_k$ with
$h=\mathrm{diag}(n)$ has energy 1 for every $k$ and a limit of energy 0.

\paragraph{The tight-enclosure obstruction.} Formally, each link $e$ of $R$
contributes $n_e/3456$ at second order, with $n_e$ the number of omitted faces
containing it (28 links with $n_e=4$, 16 with 3 and 4 with 2), which totals
$168/3456=7/144$: formally $\omega(h_R)=(7/144)\tau^2+\cdots$, about
$4.861\times10^{-18}$ at the cap. The upper end exceeds this by about
$2.016\times10^{11}$ and the lower end lies below it by a factor of about
16.9; the band is about $3.4\times10^{12}$ wide in ratio. A tight enclosure
needs a second-order upper bound on the unbounded $h_R$, an energy-weighted
state estimate that is not admitted, and it is recorded as an obstruction.

\subsubsection{The 2+1-dimensional recount}
\label{sec:bd2-dimension}
In the declared model on $\mathbb Z^2$ each site owns the links $(p,x)$ and
$(p,y)$, the face at base $p$ has owner set $\{p,p+e_x,p+e_y\}$, and each face
is its own interaction term, $V_X=-(\tau/3)W_f$. The AM2 counting is redone for
this geometry (Table~\ref{tab:bd2-dimension}): the maximal support is $p=3$,
three faces contain each site, the per-site sum is $J=|\tau|$ and the
creation expansion terminates at order 6, with
\begin{equation}
 \begin{gathered}
 G_3(t)=8e^{6t}(1+8t),\qquad G_3'(t)=8e^{6t}(14+48t),\\
 G_3(\tfrac1{64})=9e^{3/32}\approx9.88457,\qquad G_3'(\tfrac1{64})=118e^{3/32}\approx129.598 .
 \end{gathered}
 \label{eq:bd2-G3}
\end{equation}
The numerators $2^p(2p)^k(1+k(p+1)/p)$ for
$k=0,\ldots,6$ are 8, 112, 1056, 8640, 65664, 476928 and 3359232, and a
three-qubit fixture checks termination
($\operatorname{ad}_C^6(V)\Omega=-720|111\rangle$, $\operatorname{ad}_C^7(V)=0$).
With a directed series enclosure of $e^{3/32}$, the self-map condition
$JG_3(R)\le R$ at $R=1/64$ binds and gives the cap
\begin{equation}
 |\tau|\le\frac1{576\,e^{3/32}},\qquad\text{directed lower bound }\frac{1580747155173}{10^{15}}\approx1.5807\times10^{-3},
 \label{eq:bd2-cap}
\end{equation}
while the exclusion condition $2JG_3'(R)<1$ alone would allow
$1/(236e^{3/32})\approx3.858\times10^{-3}$. These constants are recounted, not
a theorem: no finite-volume ground-state or gap statement, no AQ chain, no
node and no dictionary is claimed in 2+1 dimensions, where $g^2$ has mass
dimension and the 3+1D dictionary does not apply.

\begin{table}[htbp]
\centering\small
\caption{The AM2 constants recounted in 2+1 dimensions (3+1D values not
reused).}
\label{tab:bd2-dimension}
\begin{tabular}{lll}
\toprule
Quantity & 2+1D (recounted) & 3+1D (AM2)\\
\midrule
star and faces per anchor & $\{0,e_x,e_y\}$, one face & $\{0,e_x,e_y,e_z\}$, 21 faces\\
terms containing a site & 3 faces & 4 stars, 49 faces\\
per-site sum $J$ & $|\tau|$ & $28|\tau|$\\
maximal support, termination order & 3, 6 & 4, 8\\
majorant $G_p(t)$ & $8e^{6t}(1+8t)$ & $16e^{8t}(1+10t)$\\
\bottomrule
\end{tabular}
\end{table}

\subsubsection{Obligations, no-transfer rows and code provenance}
\label{sec:bd2-obligations}
The review records six rows, each open or a no-transfer row with its missing
premise:
\begin{itemize}
\item \emph{area law:} analyticity of the reduced density in a complex
 $\tau$-disc uniformly in $N$ (a zero-free region of the complexified
 normalization) and a family of growing loops; no zero-free region is
 admitted;
\item \emph{a certified sign of the $\mathbb Z^3$ 1x2 mean:} a uniform remainder
 $|\omega(W_{1\times2})-(7/124416)\tau^2|\le K_4\tau^4$ with $K_4\tau^2$ below
 $7/124416$ for F1 boxes and the limit, and a third-order remainder for F2
 boxes;
\item \emph{a tight electric enclosure:} a second-order upper bound on the
 unbounded $\omega(h_R)$;
\item \emph{the site-blocked uniform 3+1D regime:} its own contract;
\item \emph{the finite-volume theorem in 2+1 dimensions:} a two-dimensional
 finite-volume prescription, the on-site domains and AM2 sections~4--6
 re-verified, then the AQ chain, node and dictionary;
\item \emph{the Einstein--QED evidence rounds:} a no-transfer row, since they
 share no equation with this SU(2) chain.
\end{itemize}

\paragraph{Code provenance.} The forward checker imports only the standard
library. The AZ2 certificate algebra is copied verbatim with attribution and
compared block by block with the pinned snapshot, and the copied code
reproduces the AZ2 gate enclosure of $\langle W_1\rangle$ exactly. The Round11
solver is read as bytes from its snapshot to bind 13 exact source strings,
and it is never imported or executed; the producer agent did not open it.
No interpreter cache is written under the Round11 or Round32 directories or
in the closure.

\paragraph{Evidence and review.} The forward checker runs 41 exact checks and
implements all 21 controls as damaging mutations, with 107 rejections; its
outputs are byte-identical under \texttt{-B} and \texttt{-B -O}. The skeptic's
pre-comparison package (64 checks, 21 controls, 85 rejected mutations with 4
positives) had its values final before the forward commit at 06:00:06 UTC and
was committed 10~min 2~s after it, at 06:10:08; its exposure was to input
file names, hashes and possibly the commit subject. Every exported value
equals, or certifiably contains, the skeptic's own. The post-comparison
checker runs 22 checks and 38 source-edit runs: 21 validator weakenings, 12
must-abort edits, and 3 silent edits, of which 1 is harmless and 2 are
export-only edits caught only by the skeptic's validator. It also rejects 19
damaged packets. The closure of 39 files verifies, with its 35 declared
inputs. The review found no blocking issue. It records the producer's
contract readings W1--W7: the single model id names only the graph;
evenness is admitted for F1 and the limit only, so a certified sign for F2
would need a third-order remainder; ``three faces per site'' counts the faces
whose owner set contains the site; the per-site sum is linear in $|\tau|$;
the tier name of $K_2'$; and the sign restriction of the Round11 solver
interface. It also keeps separate the two scopes of the 1x2 loop, the bound
for F1, F2 and the limit and the evenness for F1 and the limit only.

\subsubsection{Limitations recorded by the gate}
\label{sec:bd2-limits}
The BD2 gate records the following limitations verbatim.
\begin{gatequote}
\begin{itemize}
\item Graph scope: the Round11 two-square patch only, at D=6 and D=8 and the
declared grid; static means with no dynamical or mass-gap reading;
transfers\_to\_aq false; the graph coefficients are finite-graph values, and
7/124416 = (7/216)/24\textasciicircum{}2 is a labelled identity between the
graph and the formal Z\textasciicircum{}3 coefficient, not a transfer.
\item Z\textasciicircum{}3 scope: the zero-selected family at |tau| at most
10\textasciicircum{}-8, the fixed cover R and fixed spacing; no certified
value or sign of the Z\textasciicircum{}3 1x2 mean (the admitted bound
K\_2\textquotesingle{} tau\textasciicircum{}2 exceeds the formal term by about
2.38479e8; a certified sign needs a uniform fourth-order remainder, and for F2
boxes a third-order one); evenness in tau is not claimed for F2 boxes.
\item The band is two-sided but loose (ratio about 3.4e12) and is not a tight
enclosure, which would need a second-order upper bound for the unbounded h\_R;
its lower end rests on the AY1 forward F11-F14 per-box ball (a reviewed step,
not a gate sentence) for boxes and on AY2 item (2) for the limit; its upper
end passes to the limit by lower semicontinuity only.
\item No area law and no string-tension statement is made: no zero-free region
of the reduced density is admitted.
\item 2+1 dimensions: recounted AM2 constants only; no finite-volume
ground-state or gap theorem and no AQ chain, node or dictionary in 2+1
dimensions (g\textasciicircum{}2 has mass dimension); each needs its own
contract.
\item The enclosures use the admitted AZ2 certificate code copied verbatim
(Weyl gap 3-2|l|, Temple, the Round11 Schur-complement tail comparison, Eckart
and Davis-Kahan angles); the skeptic re-evaluated the lemma in its own basis
and has its own sharper tail certificate (worst D=8 relative width 3.56e-18);
Kato, Weyl, Temple, Eckart, Davis-Kahan, Peter-Weyl and Fuchs-van de Graaf are
cited, not machine-checked.
\item Review standing: model-agent skeptic with correlated ancestry (same
model family as the producer and the advisor); its pre-comparison package was
committed after the producer packet, with its values final before the producer
commit; not human peer review; scientific priority is not verified; weak
coupling and the continuum are not addressed, and no estimate is made
uniformly in the lattice spacing a.
\end{itemize}
\end{gatequote}
Full derivations and reviewed implementations:
\repo{research/round33/forward/bd2/report.md}{BD2 forward},
\repo{research/round33/skeptic/bd2.md}{independent model-agent review} and the
\repo{research/round33/contracts/bd2.json}{frozen contract}; the graph
certificate is that of the \repo{research/round32/advisor/az2-gate.json}{AZ2
gate}, which is not edited.

\subsection*{What transferred and what did not}
\label{sec:apps-summary}
Table~\ref{tab:apps-summary} collects the outcome of the applications stage.
Its rows are the admitted equations that BD1 and BD2 applied, and its columns
the models they were applied to. ``Source'' marks the model in which an
equation was admitted; every other cell reads transfer, obstruction or not
attempted, with its reason. A transfer is a labelled statement about the
named model with its own exact checks, never a statement about another model,
and ``not attempted'' records no failure. The Einstein--QED evidence rounds
share no equation with this chain and are a no-transfer row of BD2; they are
not a column.

\begin{table}[htbp]
\centering\scriptsize
\caption{What transferred and what did not. Columns: the SU(2) zero-selected
family (boxes and the limit of the named constructions); the one-plaquette
and box models $H_{FG}(G)$, $H^G_N$ of the other groups; the Round11
two-plaquette graph with independent couplings; the declared 2+1D SU(2)
model.}
\label{tab:apps-summary}
\begin{tabularx}{\linewidth}{@{}>{\raggedright\arraybackslash}p{0.14\linewidth}
 *{6}{>{\raggedright\arraybackslash}X}@{}}
\toprule
Equation & SU(2) family & SU(4), U(1), Z2 & SU(5) & SU(3), SO(3) & Two-plaquette graph & 2+1D model\\
\midrule
Link-flip lemma (AW1) & Source; used for the area parity & Transfer: central element acting as $-1$; box by box & Obstruction: $c_4=1/(2^{30}3^{10})$; no central $-1$ & Obstruction: $c_2=1/589824$, $1/331776$; no central $-1$ & Transfer: one-link flips give the parity in each coupling & Not attempted: constants recounted only\\[3pt]
First-order parity theorem (AW1) & Source (finite boxes) & Transfer: $\mathbb E[W^3]=0$ & Transfer: $\mathbb E[W^3]=0$ by the $\mathbb Z_5$ grading & Obstruction: $d\omega(W^2)/d\tau=1/6912$, $1/2592$ & Not attempted: parities come from the flips & Not attempted: no finite-volume theorem\\[3pt]
First-order Wilson coefficient $\tau/144$ (AW1) & Source; reproduced as a convention check & Transfer of the method, own values $1/2880$, $1/96$, $1/48$; $1/144$ not transferred & Transfer of the method, own value $1/5760$ & Transfer of the method, own values $1/1152$, $1/864$ & Transfer: $\langle W_1\rangle=l_1/6+\cdots$ with independent couplings & Not attempted\\[3pt]
Area parity (AW1 flip with BB2) & Transfer at $\kappa=0$: every open centered whole-star (F1) box and pointwise for the limit; not claimed for F2 boxes (so the 1x2 loop is even for F1 and the limit, not F2) & Not attempted: no limit and no regime for all $N$ & Not attempted: no flip & Not attempted: no flip & Not attempted & Not attempted\\[3pt]
AM2, AV1, AQ1--AQ2, BB1--BB2 chain & Source & Not attempted: group-specific AM2 re-instantiation is an obligation & Not attempted: obligation & Not attempted: obligation & Not attempted: finite graph, no AQ transfer & Not attempted: AM2 constants recounted, cap about $1.5807\times10^{-3}$; no theorem\\[3pt]
Coupling dictionary $\tau=96/g^4$ & Source & Not attempted: a group-specific dictionary is needed & Not attempted & Not attempted & Not attempted & Not attempted: $g^2$ has mass dimension\\[3pt]
$R$-local bound $K_2'\tau^2$ (AY1) & Transfer: $|\omega(W_{1\times2})|\le K_2'\tau^2$ for F1, F2 and the limit; certified sign an obstruction & Not attempted & Not attempted & Not attempted & Not attempted: $7/124416=(7/216)/24^2$ is a labelled identity only & Not attempted\\[3pt]
Reset budget and product distance (AQ1, AY1, AY2) & Transfer: $\tfrac32L^2\le\omega(h_R)\le98|\tau|$; tight enclosure an obstruction & Not attempted & Not attempted & Not attempted & Not attempted & Not attempted\\[3pt]
Two-plaquette certificate (AZ2) & Not attempted: finite-graph certificate, no AQ transfer & Not attempted & Not attempted & Not attempted & Transfer: independent-coupling tables through total order 4; $\langle z\rangle$ certified positive only at $l_1=l_2=\pm10^{-k}$, $k=1,2,3$ ($D=6$, $D=8$), $D=8$ relative width $\le1.79645\times10^{-17}$ & Not attempted\\
\bottomrule
\end{tabularx}
\end{table}
